In this section, we consider a celebrated DAC algorithm due to Stockmeyer~\cite{stockmeyer}, and analyze its audit complexity. Towards this end, we
first consider a QBF formula that plays a significant role in Stockmeyer's algorithm.   For a given Boolean formula $F$ on $n$ variables, and for $1 \le m \le n$, consider
  \begin{align*}
 	\varphi_{stock}^F(m):=  \exists h_1\ldots \exists h_m, \forall \vecbold{\vecbold{z_1}}, \forall \vecbold{\vecbold{z_2}}~~~~~~~~~~~~~~~~~~~~~~~~~~~~~~~~~~~~~\\ \bigvee_{i=1}^m \big( F(\vecbold{\vecbold{z_1}})\land (h_i(\vecbold{\vecbold{z_1}})=h_i(\vecbold{\vecbold{z_2}}))\rightarrow \neg F(\vecbold{\vecbold{z_2}})\big)
 \end{align*}
 
In the above formula, the $h_i$'s are hash functions from
$\mathcal{H}(n, m, 2)$.  The formula says that there exist $m$ such hash
functions such that for any two assignments
$\vecbold{z_1}, \vecbold{z_2}\in \{0,1\}^n$, if $\vecbold{z_1}$ is a
solution for $F$ and some hash function $h_i$ maps $\vecbold{z_1}$ and
$\vecbold{z_2}$ to the same cell, then $\vecbold{z_2}$ cannot be a
solution. In other words, for each solution some hash function puts it
in a cell where it is the only solution. {Recall that $\exists h_i$ is
shorthand for existentially quantifying over coefficients of $h$.}

Now, the following claim and lemma provide necessary and sufficient conditions for $\varphi_{stock}^F(m)$ evaluating to True. 


\begin{claim}[Necessary]
 	\label{cl:stockup}
 	If $\varphi^F_{stock}(m)=1$, then  $\sol{F}\leq m\cdot 2^m$.

\end{claim}
\fullversion{
\begin{proof} To see this, note that if $\varphi_{stock}^F(m)=1$, there
  	exists an injective map from $\sol{F}$ to $[m] \times \{0,1\}^m$.
  	The solution $z$ of $F$ can be uniquely identified a pair $(i,h_i(z))$ where $h_i$ is the hash function that puts it in a bucket $h_i(z)$. Hence, we get $\sol{F}\leq m\cdot 2^m$.
  \end{proof}
 }
 
 

The following lemma capturing sufficient conditions is inherent in Sipser's work~\cite{Sipser83}; to obtain the precise constant for our analysis, we provide an alternate simpler argument \shortversion{(proofs can be found in the full version at~\cite{audit-arxiv})}.
\begin{lemma}[Sufficient]
	\label{lm:stockexist}
If $\sol{F}\leq 2^{m-2}$, then $\varphi_{stock}^F(m) =1$ 
\end{lemma}






\fullversion{
\begin{proof}
For 
$h_i \xleftarrow{R} H(n,m,2)$ and a fixed $y \in \satisfying{F}$, let us define indicator variable $\gamma_{z,y,h_i}$ for all $z \in \satisfying{F} \setminus \{y\}$ as
\begin{align*}
	\gamma_{z,y,h_i} = 
	\begin{cases}
		1 & h_i(z) = h_i(y)\\
		0 & \text{otherwise}
	\end{cases}
\end{align*} 

Since $h_i$ is chosen from pairwise independent 
hash family, we have $\Pr[\gamma_{z,y,h_i} = 1] \leq \frac{1}{2^m}$.  Let $\Gamma_{F,y,h_i} = \sum_{z\in \satisfying{F}\setminus \{y\}} \gamma_{z,y,h_i}$. We have $\expect[\Gamma_{F,y,h_i}] = \sum\limits_{z	\in \satisfying{F} \setminus \{y\}} \expect[\gamma_{z,y,h_i}] = \frac{\sol{F}-1}{2^m} \leq \frac{1}{4}$, since $\sol{F}\leq 2^{m-2}$.

Now, by Markov's inequality, we have 
\begin{align*}
	\Pr[\Gamma_{F,y,h_i} \geq 1 ] \leq \expect[\Gamma_{F,y,h_i}] \leq \frac{1}{4}
\end{align*}

Therefore, we have 
\begin{align*}
	\Pr\left[ \bigwedge_{i=1}^m  \{\Gamma_{F,y,h_i} \geq 1\} \right] \leq \frac{1}{4^m}	
\end{align*} 

Now, taking union bound, we have 
\begin{align*}
	\Pr\left[ \exists y \in \satisfying{F},  \bigwedge_{i=1}^m  \{\Gamma_{F,y,h_i} \geq 1\} \right] \leq \frac{2^{m-2}}{4^m} \leq \frac{1}{2^m}	
\end{align*} 
Taking the complement of the above statement, we have 
\begin{align*}
	\Pr\left[ \forall y \in \satisfying{F},  \bigvee_{i=1}^m  \{\Gamma_{F,y,h_i} < 1\} \right] \geq 1- \frac{1}{2^m} > 0	
\end{align*}

Therefore, we have $\varphi_{stock}^F(m) =1$.
\end{proof}
}
We can now write Stockmeyer's algorithm as follows:

\begin{algorithm}
	\caption{Stock$(F)$}
		\label{algo:stock}
	\begin{algorithmic}[1]
		\State $\val \gets 0$; $\mathsf{hashassgn}_{stock}\gets \emptyset,\Cest\gets 0$;
		\State $F' \gets \mathsf{MakeCopies}(F,\log n)$ \hfill / {$F'$ has $n\log n $ variables}
		\For{m=1 to $n\log n$}
		\State  $(\mathsf{ret},\mathsf{assign})\gets \twoqbfcheck (\varphi_{stock}^{F'}(m)$);\label{line:qbf-call}
		\If{$\mathsf{ret}==1$
		}\label{line:equalcells-check}
		\State $\val=m$;
		\State $\mathsf{hashassgn}_{stock}=\mathsf{assign}$;
		\EndIf
		\State \textbf{break}
		\EndFor
		\State $\Cest=2^{(\frac{\val}{\log (n)})}$;\\
		\Return $(\val, \Cest,\hashstock)$ 
	\end{algorithmic}

\end{algorithm}


Equipped with necessary and sufficient conditions for the satisfiability of $\sol{F}\leq 2^{m-2}$, we now describe Stockmeyer's algorithm, as presented in Algorithm~\ref{algo:stock}.
We first invoke  
$\mathsf{MakeCopies}(F,\log n)$ subroutine, which makes $\log(n)$ many copies of $F$, each with a fresh set of variables and conjuncts them.
Now, the algorithm runs over $m$ from $1$ to number of variables of $F'$, i.e., $n\log n$ and checks if $\varphi_{stock}^{F'}(m)$ evaluates to 1. Each such call is a \twoqbfcheck\  as $\varphi_{stock}^{F'}$ is a $\Sigma_2^P$ formula.

Now, if $\varphi_{stock}^{F'}(i)=1$, then $\varphi_{stock}^{F'}(i+1)=1$ as well, by the definition of the formula, so the above algorithm essentially searches for the first point (smallest value of $m$) where $\varphi_{stock}^{F'}(i)=1$ and returns the estimate of number of solutions at that point as $2^{\frac{\val}{\log n}}$. Thus, when algorithm breaks at line 8, we have $\varphi_{stock}^{F'}(\val)=1$ and $\varphi_{stock}^{F'}(\val-1)=0$.

 
\begin{theorem}[Stockmeyer]
	Algorithm~\ref{algo:stock} makes O($n \log n$) queries to $\Sigma_2^P$ oracle and solves DAC. 
\end{theorem}

\fullversion{
\begin{proof}
	
From Claim~\ref{cl:stockup}, we have $\sol{F'}\leq \val\cdot 2^\val$. Taking the contrapositive of Lemma~\ref{lm:stockexist} at $\val-1$, we get $\sol{F'}>2^{(\val-1)-2}$. Combining the two, we have $2^{\val-3}\leq  \sol{F'}\leq \val \cdot 2^\val$, i.e., an $8v$-factor approximation.

Now, we can see that this results in a constant-factor approximation for $\sol{F}$, due to having taken $\log(n)$ copies of $F$. First note that $\sol{F'}=(\sol{F})^{\log(n)}$. Second we note that $\val\leq n \log n\leq n^2$ and so $\frac{1}{n^2}\leq \frac{1}{v}$. Now, by above argument we have,

\begin{align*}
  2^{\val-3}\leq \sol{F'}\leq \val\cdot 2^\val\\
  \frac{\sol{F'}}{n^2}\leq  \frac{\sol{F'}}{\val}\leq 2^\val\leq \sol{F'}\cdot 2^{3}\\
       {\left(\frac{\sol{F}}{16}\right)}^{\log(n)}\leq (2^{\frac{\val}{\log(n)}})^{\log(n)}\leq {(\sol{F})}^{\log(n)} \cdot 2^{3}\\
       \frac{\sol{F}}{16} \leq \Cest\leq 16\cdot \sol{F} \ \text{for $n\geq 4$}

\end{align*}

Thus, Algorithm~\ref{algo:stock}  computes a 16-factor approximation of $\sol{F}$ for $n \ge 4$, and makes $\bigO\big(n \log n\big)$ calls to a $\Sigma_2^P$-oracle.
\end{proof}
}
\subsection{Auditing Stockmeyer's algorithm}

Now to check the result/count given by the above algorithm, what can we do? In other words, what could be considered as a {\em certificate} for the answer? If we inspect the above algorithm, we can see that the crux is to identify the smallest value of $m$ for which 2QBFCheck returns True. 

Thus, at $\val$, we need to check that the True answer is correct, i.e., the $\val^{th}$ call to the $\Sigma_2^P$ oracle. This is easy since we can use the certificate returned by the call, i.e., the set of hash functions $h_1, \ldots, h_\val$ that was used, i.e., $\hashstock$. Given this certificate (indeed a solver returns this) we can pass it to the verifier and hence, a query to NP oracle suffices. 

However, this does not suffice. We also need to check that at $\val-1$, the False answer is correct, i.e., at $\val-1$, we need to check that $\varphi_{stock}^{F'}(\val-1)=0$, which can be answered by a query to $\Sigma_2^P$ oracle. More precisely, for $m=\val-1$, we need to check that $\neg \varphi^{F'}_{stock}(\val-1)=1$, i.e., $\forall h_1\ldots \forall h_m$, $\exists \vecbold{z_1}, \exists \vecbold{z_2} \bigwedge_{i=1}^m (F(\vecbold{z_1})\land (h_i(\vecbold{z_1})=h_i(\vecbold{z_2}))\land F(\vecbold{z_2}))$. 


\begin{algorithm}
	\caption{StockAudit$(F,\mathsf{Stock}(F))$}
	\begin{algorithmic}[1]
\State {$\mathsf{poscheck}\gets 0$; $\mathsf{negcheck}\gets 0$;}
\State  $\mathsf{poscheck}\gets \conpcheck (\varphi_{stock}^{F'}(\val)[h_1\ldots,h_v \gets \hashstock])$; \label{line:conpcheck}
                \State  $\mathsf{negcheck}\gets \twoqbfcheck(\neg\varphi_{stock}^{F'}(\val-1))$;\label{line:twoqbfcheck}   
		\If{$\mathsf{poscheck}\land \mathsf{negcheck}==1$}\label{line:equalcheck}                            
                \State \Return $\mathsf{Verified}$;
\EndIf
\end{algorithmic}
        \label{algo:stock-audit}
\end{algorithm}


Thus, the complexity is dominated by the ``False'' check at $\val-1$. In fact combining the two checks as done in line 4, we just require a single call to 2-QBF solver. Thus the \emph{audit-complexity}, or the number of variables on which this call is made is $n^2\log^2(n)+ 2 n\log n$. 

Thus we obtain:

\begin{corollary}
\label{cor:stockaudit}
The audit complexity of Algorithm~\ref{algo:stock} is $O(n^2 \log^2 n)$.

\end{corollary}





\fullversion{
\begin{proof}
  Algorithm~\ref{algo:stock-audit} gives the audit algorithm for
  Algorithm~\ref{algo:stock}.  The two checks in lines $2$ and $3$ of
  Algorithm~\ref{algo:stock-audit} can be combined into one
  check as follows.

  Let $\varphi_{st\_aud}^{F'}(v)$ be defined as

  \begin{align*}
  \forall \boldsymbol{z_1} \forall \boldsymbol{z_2} \bigvee_{i=1}^v \big(F'(\boldsymbol{z_1}) \wedge (h_i(\boldsymbol{z_1}) = h_i(\boldsymbol{z_2})) \rightarrow \neg F'(\boldsymbol{z_2})\big)[h_1, \ldots h_v \gets \hashstock]\\
  
 \bigwedge \forall h_1 \ldots \forall h_{v-1} \exists \boldsymbol{z_1} \exists \boldsymbol{z_2} \bigwedge_{i=1}^{v-1} \big(F'(\boldsymbol{z_1}) \wedge (h_i(\boldsymbol{z_1}) = h_i(\boldsymbol{z_2})) \wedge F'(\boldsymbol{z_2})\big)~~~~~~~~~~~~~~~~~~~~~~~~~~~~
  \end{align*}
\normalsize
Recalling that $\forall h_i$ is simply syntactic sugar for quantifying
over coefficients of the hash functions, the above formula is a formula
with quantifier prefix $\forall^* \exists^*$ and its truth can be
decided by one invocation of a $\Sigma_2^P$ oracle.

To determine the audit complexity, notice that each $\boldsymbol{z_i}$
in the above formula represents $n \log n$ bits, since $F'$ has a
support of size $n \log n$.  Similarly, each $h_i$ is a hash function
from $\mathcal{H}(n\log n, v-1, 2)$, where $v \le n \log n$. Since
$v \le n \log n$, by our discussion in Section~\ref{sec:prelim}, each
such hash function can be represented using $\bigO(n \log n)$ bits.
Therefore, the total number of Boolean variables in the quantifier
prefix of $\varphi_{staud}^{F'}(v)$ is in $\bigO\big( (n \log n)^2 +
n\log n\big) = \bigO(n^2 \log^2 n)$.







\end{proof}
}
